\section{USAMO 2011}
        \begin{enumerate}
        	\item (\textit{USAMO 2011}) Let $a$, $b$, $c$ be positive real numbers such that $a^2 + b^2 + c^2 + (a + b + c)^2 \le 4$.  Prove that

 \[\frac{ab + 1}{(a + b)^2} + \frac{bc + 1}{(b + c)^2} + \frac{ca + 1}{(c + a)^2} \ge 3.\]


 



\textbf{Solution 1}

Since
 \begin{align*} (a+b)^2 + (b+c)^2 + (c+a)^2 &= 2(a^2 + b^2 + c^2 + ab + bc + ca) \\ 	&= a^2 + b^2 + c^2 + (a + b + c)^2, \end{align*}
it is natural to consider a change of variables:
 \begin{align*} \alpha &= b + c \\ \beta &= c + a \\ \gamma &= a + b \end{align*}
with the inverse mapping given by:
 \begin{align*} a &= \frac{\beta + \gamma - \alpha}2 \\ b &= \frac{\alpha + \gamma - \beta}2 \\ c &= \frac{\alpha + \beta - \gamma}2 \end{align*}
With this change of variables, the constraint becomes
 \[\alpha^2 + \beta^2 + \gamma^2 \le 4,\]while the left side of the inequality we need to prove is now
 \begin{align*} & \frac{\gamma^2 - (\alpha - \beta)^2 + 4}{4\gamma^2} + \frac{\alpha^2 - (\beta - \gamma)^2 + 4}{4\alpha^2} + \frac{\beta^2 - (\gamma - \alpha)^2 + 4}{4\beta^2} \ge \\ & \frac{\gamma^2 - (\alpha - \beta)^2 + \alpha^2 + \beta^2 + \gamma^2}{4\gamma^2} + \frac{\alpha^2 - (\beta - \gamma)^2 + \alpha^2 + \beta^2 + \gamma^2}{4\alpha^2} + \frac{\beta^2 - (\gamma - \alpha)^2 + \alpha^2 + \beta^2 + \gamma^2}{4\beta^2} = \\ & \frac{2\gamma^2 + 2\alpha\beta}{4\gamma^2} + \frac{2\alpha^2 + 2\beta\gamma}{4\alpha^2} + \frac{2\beta^2 + 2\gamma\alpha}{4\beta^2} = \\ & \frac32 + \frac{\alpha\beta}{2\gamma^2} + \frac{\beta\gamma}{2\alpha^2} + \frac{\gamma\alpha}{2\beta^2}. \end{align*}
Therefore it remains to prove that
 \[\frac{\alpha\beta}{2\gamma^2} + \frac{\beta\gamma}{2\alpha^2} + \frac{\gamma\alpha}{2\beta^2} \ge \frac32.\]We note that the product of the three (positive) terms is 1/8, therefore by AM-GM their mean is at least 1/2, and thus their sum is at least 3/2 and we are done.

 



\textbf{Solution 2}

Rearranging the condition yields that
 \[a^2 + b^2 + c^2 +ab+bc+ac \le 2\]

 Now note that
 \[\frac{2ab+2}{(a+b)^2} \ge \frac{2ab+a^2 + b^2 + c^2 +ab+bc+ac}{(a+b)^2}=\frac{(a+b)^2 + (c+a)(c+b)}{(a+b)^2}\]

 Summing this for all pairs of $\{ a,b,c \}$ gives that
 \[\sum_{cyc} \frac{2ab+2}{(a+b)^2} \ge 3+ \sum_{cyc}\frac{(c+a)(c+b)}{(a+b)^2} \ge 6\]

 By AM-GM. Dividing by $2$ gives the desired inequality.

 



\textbf{Solution 3}

Given $a^2+b^2+c^2+(a+b+c)^2 \leq 4$, prove that
 \[\frac{ab+1}{(a+b)^2}+\frac{bc+1}{(b+c)^2}+\frac{ac+1}{(a+c)^2} \geq 3.\]

 From the given constraint $a^2+b^2+c^2+(a+b+c)^2 \leq 4$, we can expand to obtain:
 \[a^2+b^2+c^2+a^2+b^2+c^2+2ab+2bc+2ca \leq 4\]
 \[2(a^2+b^2+c^2)+2(ab+bc+ca) \leq 4\]
 \[\frac{a^2+b^2+c^2}{2}+\frac{ab+bc+ca}{2} \leq 1.\]

 Now, observe that we can rewrite the left-hand side of our inequality as:\begin{align*}&\sum_{\text{cyc}} \frac{ab+1}{(a+b)^2} \\&= \sum_{\text{cyc}} \frac{ab+\frac{a^2+b^2+c^2}{2}+\frac{ab+bc+ca}{2}}{(a+b)^2} \\&= \sum_{\text{cyc}} \left[\frac{ab+\frac{a^2+b^2}{2}}{(a+b)^2} + \frac{\frac{c^2}{2}+\frac{bc+ca}{2}}{(a+b)^2}\right] \\&= \sum_{\text{cyc}} \left[\frac{ab+\frac{a^2+b^2}{2}}{(a+b)^2}\right] + \sum_{\text{cyc}} \left[\frac{\frac{a^2+c^2}{2}}{(a+c)^2} + \frac{\frac{ac+bc}{2}}{(a+b)^2}\right].\end{align*}


 We now evaluate each part separately.Note that 
 \[\frac{ab+\frac{a^2+b^2}{2}}{(a+b)^2} = \frac{\frac{(a+b)^2}{2}}{(a+b)^2} = \frac{1}{2}.\]Therefore, 
 \[\sum_{\text{cyc}} \frac{ab+\frac{a^2+b^2}{2}}{(a+b)^2} = \frac{3}{2}.\]By the RMS-AM inequality, we have $(a+c)^2 \leq 2(a^2+c^2)$, which gives $\frac{1}{(a+c)^2} \geq \frac{1}{2(a^2+c^2)}$. Thus,\begin{align*}\sum_{\text{cyc}} \frac{\frac{a^2+c^2}{2}}{(a+c)^2} &\geq \sum_{\text{cyc}} \frac{\frac{a^2+c^2}{2}}{2(a^2+c^2)} \\&= \sum_{\text{cyc}} \frac{1}{4} = \frac{3}{4}.\end{align*}
The remaining terms can be written as:\begin{align*}\sum_{\text{cyc}} \frac{\frac{ac+bc}{2}}{(a+b)^2} &= \frac{1}{2}\sum_{\text{cyc}} \frac{c(a+b)}{(a+b)^2} \\&= \frac{1}{2}\sum_{\text{cyc}} \frac{c}{a+b} \\&= \frac{1}{2}\left(\frac{a}{b+c}+\frac{b}{a+c}+\frac{c}{a+b}\right).\end{align*}
By Nesbitt's inequality, 
 \[\frac{a}{b+c}+\frac{b}{a+c}+\frac{c}{a+b} \geq \frac{3}{2},\]so this part contributes at least $\frac{3}{4}$.Adding all three parts together:
 \[\frac{ab+1}{(a+b)^2}+\frac{bc+1}{(b+c)^2}+\frac{ac+1}{(a+c)^2} \geq \frac{3}{2}+\frac{3}{4}+\frac{3}{4} = 3.\]

 Therefore, the inequality is proven. $\square$

 


	\item (\textit{USAMO 2011}) An integer is assigned to each vertex of a regular pentagon so that the sum of the five integers is 2011. A turn of a solitaire game consists of subtracting an integer $m$ from each of the integers at two neighboring vertices and adding $2m$ to the opposite vertex, which is not adjacent to either of the first two vertices. (The amount $m$ and the vertices chosen can vary from turn to turn.) The game is won at a certain vertex if, after some number of turns, that vertex has the number 2011 and the other four vertices have the number 0. Prove that for any choice of the initial integers, there is exactly one vertex at which the game can be won.


 



\textbf{Solution}

Let $\mathbb{F}_5$ be the field of positive residues modulo 5. We label the vertices of the pentagon clockwise with the residues 0, 1, 2, 3 and 4.For each $i \in \mathbb{F}_5$ let $n_i$ be the integer at vertex $i$ and let $r_i \in \mathbb{F}_5$ be defined as:
 \[r_i   \equiv 4n_{i+1} + 3n_{i+2} + 2n_{i+3} + n_{i+4}   \equiv \sum_{k\in\mathbb{F}_5}(i-k)n_k \pmod 5.\]Let $s = \sum_{i\in\mathbb{F}_5} n_i$. A move in the game consists of
 \[(n_i, n_{i+1}, n_{i+2}, n_{i+3}, n_{i+4}) \mapsto (n_i + 2m, n_{i+1}, n_{i+2} - m, n_{i+3} - m, n_{i+4})\]for some vertex $i \in \mathbb{F}_5$ and integer $m$. We immediately see that $s$ is an invariant of the game. After our move the new value of $r_i$ is decreased by $3m + 2m \equiv 0 \pmod 5$ as a result of the change in the $3n_{i+2}$ and $2n_{i+3}$ terms. So $r_i$ does not change after a move at vertex $i$.

 For all $i \in \mathbb{F}_5$ we have:
 \[r_{i+1} - r_i   \equiv \sum_{k\in\mathbb{F}_5} ((i+1-k)n_k - (i-k)n_k)   \equiv \sum_{k\in\mathbb{F}_5} n_k \equiv s \pmod 5.\]

 Therefore, the $r_i$ form an arithmetic progression in $\mathbb{F}_5$ with a difference of $s$. Since $r_k$ is unchanged by a move at vertex $k$, so are all the remaining $r_i$ as the differences are constant.

 Provided $s \not\equiv 0 \pmod 5$, we see that the mapping $i \mapsto r_i$ is a bijection $\mathbb{F}_5 \to \mathbb{F}_5$ and exactly one vertex will have $r_i \equiv 0 \pmod 5$. As $r_i$ is an invariant, a winning vertex must have $r_i \equiv 0$, since in the final state each $n_k$ with $k \ne i$ is zero. So, for $s \not\equiv 0 \pmod 5$, if a winning vertex exists, it is the unique vertex with $r_i \equiv 0$.

 Without loss of generality, it remains to show that if $r_0 \equiv 0 \pmod 5$, then 0 must be a winning vertex. To prove this, we perform the following moves:
 \begin{align*} (n_0, n_1, n_2, n_3, n_4) &\mapsto (n_0-n_1, 0, n_2, n_3 + 2n_1, n_4), & (i, m) &= (3, n_1) \\  &\mapsto (n_0-n_1-n_4, 0, n_2+2n_4, n_3 +2n_1, 0) & (i, m) &= (2, n_4) \end{align*}
We designate the new state $(n'_0, 0, n'_2, n'_3, 0)$. Since $r_0 \equiv 0$ is an invariant, and $n'_1 = n'_4 = 0$, we now have $r_0 \equiv n'_3 - n'_2 = 5p$, for some integer $p$. Our final set of moves is:
 \begin{align*} (n'_0, 0, n'_2, n'_3, 0) &\mapsto (n'_0+p, p, n'_2, n'_3 - 2p, 0), & (i, m) &= (3, -p)\\  &\mapsto (n'_0+p, 0, n'_2-p, n'_3 - 2p, 2p), & (i, m) &= (4, p) \\  &\mapsto (n'_0-p, 0, n'_2+3p, n'_3 - 2p, 0), & (i, m) &= (2, 2p)\\  &\mapsto (n'_0+2n'_2 + 5p, 0, 0, n'_3-n'_2 - 5p = 0, 0) & (i, m) &= (0, n'_2 + 3p) \end{align*}
Now our chosen vertex 0 is the only vertex with a non-zero value, and since $s$ is invariant, that value is $s$ as required. Since a vertex $i$ with $r_i \equiv 0 \pmod 5$ is winnable, and with $s = 2011 \not\equiv 0 \pmod 5$ we always have a unique such vertex, we are done.

 


	\item (\textit{USAMO 2011}) In hexagon $ABCDEF$, which is nonconvex but not self-intersecting, no pair of opposite sides are parallel.  The internal angles satisfy $\angle A = 3 \angle D$, $\angle C = 3 \angle F$, and $\angle E = 3 \angle B$.  Furthermore, $AB = DE$, $BC = EF$, and $CD = FA$.  Prove that diagonals $\overline{AD}$, $\overline{BE}$, and $\overline{CF}$ are concurrent.


 



\textbf{Solution 1}

Let $\angle D = \alpha$, $\angle F = \gamma$, and $\angle B = \beta$, $AB=DE=p$, $BC=EF=q$, $CD=FA=r$.  Define the vectors: 
 \[\vec{u} = \vec{AB} + \vec{DE}\] 
 \[\vec{v} = \vec{BC} + \vec{EF}\] 
 \[\vec{w} = \vec{CD} + \vec{FA}\] Clearly, $\vec{u}+\vec{v}+\vec{w}=\textbf{0}$.

 Let $AB$ intersect $DE$ at $X$. Note that $\angle X = 360^\circ - \angle D - \angle C - \angle B = 360^\circ - \alpha - 3\gamma - \beta = 180^\circ - 2\gamma$. Define the points $M$ and $N$ on lines $AB$ and $DE$ respectively so that $\vec{MX} = \vec{AB}$ and $\vec{XN} = \vec{DE}$. Then $\vec{u} = \vec{MN}$.  As $XMN$ is isosceles with $XM = XN = p$, the base angles are both $\gamma$.  Thus, $|\vec{u}|=2p \cos \gamma$.  Similarly, $|\vec{v}|=2q \cos \alpha$ and $|\vec{w}| = 2r \cos \beta$.

 Next we will find the angles between $\vec{u}$, $\vec{v}$, and $\vec{w}$.  As $\angle MNX = \gamma$, the angle between the vectors $\vec{u}$ and $\vec{NE}$ is $\gamma$.  Similarly, the angle between $\vec{EF}$ and $\vec{v}$ is $\alpha$. Since the angle between $\vec{NE}$ and $\vec{EF}$ is $\angle E = 3\beta$, the angle between $\vec{u}$ and $\vec{v}$ is $360^\circ - \gamma - 3\beta - \alpha = 180^\circ - 2\beta$.  Similarly, the angle between $\vec{v}$ and $\vec{w}$ is $180^\circ - 2\gamma$, and the angle between $\vec{w}$ and $\vec{u}$ is $180^\circ - 2\alpha$.

 And since $\vec{u}+\vec{v}+\vec{w}=\vec{0}$, we can arrange the three vectors to form a triangle, so the triangle with sides of lengths $2p \cos \gamma$, $2q \cos \alpha$, and $2r \cos \beta$ has opposite angles of $180^\circ - 2\gamma$, $180^\circ - 2\alpha$, and $180^\circ - 2\beta$, respectively.  So by the law of sines: 
 \[\frac{2p \cos \gamma}{\sin 2\gamma} = \frac{2q \cos \alpha}{\sin 2\alpha} = \frac{2r \cos \beta}{\sin 2\beta}\] 
 \[\frac{p}{\sin \gamma} = \frac{q}{\sin \alpha} = \frac{r}{\sin \beta},\] and the triangle with sides of length $p$, $q$, and $r$ has corresponding angles of $\gamma$, $\alpha$, and $\beta$. It follows by SAS congruency that this triangle is congruent to $FAB$, $BCD$, and $DEF$, so $FD=p$, $BF=q$, and $BD=r$, and $D$, $F$, and $B$ are the reflections of the vertices of triangle $ACE$ about the sides.  So $AD$, $BE$, and $CF$ concur at the orthocenter of triangle $ACE$.

 



\textbf{Solution 2}

We work in the complex plane, where lowercase letters denote their corresponding point's poition. Let $P$ denote hexagon $ABCDEF$. Since $AB=DE$, the condition $AB\not\parallel DE$ is equivalent to $a-b+d-e\ne 0$.

 Construct a "phantom hexagon" $P'=A'B'C'D'E'F'$ as follows: let $A'C'E'$ be a triangle with $\angle{A'C'E'}=\angle{F}$, $\angle{C'E'A'}=\angle{B}$, and $\angle{E'A'C'}=\angle{F}$ (this is possible since $\angle{B}+\angle{D}+\angle{F}=180^\circ$ by the angle conditions), and reflect $A',C',E'$ over its sides to get points $D',F',B'$, respectively. By rotation and reflection if necessary, we assume $A'B'\parallel AB$ and $P',P$ have the same orientation (clockwise or counterclockwise), i.e. $\frac{b-a}{b'-a'}\in\mathbb{R}^+$. It's easy to verify that $\angle{X'}=\angle{X}$ for $X\in\{A,B,C,D,E,F\}$ and opposite sides of $P'$ have equal lengths. As the corresponding sides of $P$ and $P'$ must then be parallel, there exist positive reals $r,s,t$ such that $r=\frac{a-b}{a'-b'}=\frac{d-e}{d'-e'}$, $s=\frac{b-c}{b'-c'}=\frac{e-f}{e'-f'}$, and $t=\frac{c-d}{c'-d'}=\frac{f-a}{f'-a'}$. But then $0\ne a-b+d-e=r(a'-b'+d'-e')$, etc., so the non-parallel condition "transfers" directly from $P$ to $P'$ and
 \begin{align*} 0 &=(a-b+d-e)+(b-c+e-f)+(c-d+f-a) \\ &=r(a'-b'+d'-e')+s(b'-c'+e'-f')+t(c'-d'+f'-a') \\ &=(r-t)(a'-b'+d'-e')+(s-t)(b'-c'+e'-f'). \end{align*}
If $r-t=s-t=0$, then $P$ must be similar to $P'$ and the conclusion is obvious.

 Otherwise, since $a'-b'+d'-e'\ne0$ and $b'-c'+e'-f'\ne0$, we must have $r-t\ne0$ and $s-t\ne0$. Now let $x=\frac{a'+d'}{2}$, $y=\frac{c'+f'}{2}$, $z=\frac{e'+b'}{2}$ be the feet of the altitudes in $\triangle{A'C'E'}$; by the non-parallel condition in $P'$, $x,y,z$ are pairwise distinct. But $\frac{z-x}{z-y}=\frac{s-t}{r-t}\in\mathbb{R}$, whence $x,y,z$ are three distinct collinear points, which is clearly impossible. (The points can only be collinear when $\triangle{A'C'E'}$ is a right triangle, but in this case two of $x,y,z$ must coincide.)

 Alternatively (for the previous paragraph), WLOG assume that $(A'C'E')$ is the unit circle, and use the fact that $b'=a'+c'-\frac{a'c'}{e'}$, etc. to get simple expressions for $a'-b'+d'-e'$ and $b'-c'+e'-f'$.

 



\textbf{Solution 3}

We work in the complex plane to give (essentially) a complete characterization when the parallel condition is relaxed.

 WLOG assume $a,b,c$ are on the unit circle. It suffices to show that $a,b,c$ uniquely determine $d,e,f$, since we know that if we let $E$ be the reflection of $B$ over $AC$, $D$ be the reflection of $A$ over $CE$, and $F$ be the reflection of $C$ over $AE$, then $ABCDEF$ satisfies the problem conditions. (*)

 It's easy to see with the given conditions that
 \begin{align*} (a-b)(c-d)(e-f) &= (b-c)(d-e)(f-a) \Longleftrightarrow f=\frac{(a-b)(c-d)e+(c-b)(e-d)a}{(a-b)(c-d)+(c-b)(e-d)} \\ \frac{(e-a)(c-b)}{(a-b)(c-d)+(c-b)(e-d)} = \frac{f-e}{d-e} &= \left(\frac{c-b}{a-b}\right)^2 \overline{\left(\frac{a-b}{c-b}\right)} = \frac{c-b}{a-b}\cdot\frac{c}{a} \Longleftrightarrow d=\frac{c[(a-b)c+(c-b)e]+a(a-e)(a-b)}{c[(a-b)+(c-b)]} \\ \frac{(a-b)(c-d)+(c-b)(e-d)}{(a-e)(c-d)} = \frac{b-a}{f-a} &= \left(\frac{e-d}{c-d}\right)^2 \overline{\left(\frac{c-d}{e-d}\right)}. \end{align*}
Note that
 \[\frac{e-d}{c-d}=\frac{(a-b)[c(e-c)+a(e-a)]}{c(c-e)(c-b)-a(a-e)(a-b)},\]so plugging into the third equation we have
 \begin{align*} \frac{a(a-b)(2b-a-c)}{c(c-e)(c-b)-a(a-e)(a-b)} &=\frac{(a-b)+(c-b)\frac{(a-b)[c(e-c)+a(e-a)]}{c(c-e)(c-b)-a(a-e)(a-b)}}{(a-e)}\\ &=\left(\frac{(a-b)[c(e-c)+a(e-a)]}{c(c-e)(c-b)-a(a-e)(a-b)}\right)^2\overline{\left(\frac{c(c-e)(c-b)-a(a-e)(a-b)}{(a-b)[c(e-c)+a(e-a)]}\right)}\\ &=\left(\frac{(a-b)[c(e-c)+a(e-a)]}{c(c-e)(c-b)-a(a-e)(a-b)}\right)^2\frac{\frac{1}{c}\left(\overline{e}-\frac{1}{c}\right)\frac{b-c}{bc}-\frac{1}{a}\left(\overline{e}-\frac{1}{a}\right)\frac{b-a}{ba}}{\frac{b-a}{ab}\left(\frac{1}{c}\left(\frac{1}{c}-\overline{e}\right)+\frac{1}{a}\left(\frac{1}{a}-\overline{e}\right)\right)}\\ &=\left(\frac{(a-b)[c(e-c)+a(e-a)]}{c(c-e)(c-b)-a(a-e)(a-b)}\right)^2\frac{c^3(a\overline{e}-1)(a-b)-a^3(c\overline{e}-1)(c-b)}{c(a-b)[a^2(c\overline{e}-1)+c^2(a\overline{e}-1)]}\\ &=\frac{(a-b)[c(e-c)+a(e-a)]^2}{[c(c-e)(c-b)-a(a-e)(a-b)]^2}\frac{c^3(a\overline{e}-1)(a-b)-a^3(c\overline{e}-1)(c-b)}{c[a^2(c\overline{e}-1)+c^2(a\overline{e}-1)]}. \end{align*}
Simplifying, this becomes
 \begin{align*} &ac(2b-a-c)[c(c-e)(c-b)-a(a-e)(a-b)][a^2(c\overline{e}-1)+c^2(a\overline{e}-1)]\\ &=[c(e-c)+a(e-a)]^2[c^3(a\overline{e}-1)(a-b)-a^3(c\overline{e}-1)(c-b)]. \end{align*}
Of course, we can also "conjugate" this equation -- a nice way to do this is to note that if
 \[x=\frac{(a-b)[c(e-c)+a(e-a)]}{c(c-e)(c-b)-a(a-e)(a-b)},\]then
 \[\frac{a(2b-a-c)}{c(e-c)+a(e-a)}=\frac{x}{\overline{x}}=\overline{\left(\frac{c(e-c)+a(e-a)}{a(2b-a-c)}\right)}=\frac{\frac{1}{c}\left(\overline{e}-\frac{1}{c}\right)+\frac{1}{a}\left(\overline{e}-\frac{1}{a}\right)}{\frac{1}{a}\left(\frac{2}{b}-\frac{1}{a}-\frac{1}{c}\right)},\]whence
 \[ac(2b-a-c)[2ac-b(a+c)]=b[c(e-c)+a(e-a)][a^2(c\overline{e}-1)+c^2(a\overline{e}-1)].\]If $a+c\ne 0$, then eliminating $\overline{e}$, we get
 \[e\in\left\{a+c-\frac{ac}{b},a+\frac{2c(c-b)}{a+c},c+\frac{2a(a-b)}{a+c}\right\}.\]The first case corresponds to (*) (since $a,b,c,e$ uniquely determine $d$ and $f$), the second corresponds to $AB\parallel DE$ (or equivalently, since $AB=DE$, $a-b=e-d$), and by symmetry, the third corresponds to $CB\parallel FE$.

 Otherwise, if $c=-a$, then we easily find $b^2e=a^4\overline{e}$ from the first of the two equations in $e,\overline{e}$ (we actually don't need this, but it tells us that the locus of working $e$ is a line through the origin). It's easy to compute $d=e+\frac{a(a-b)}{b}$ and $f=e+\frac{a(a+b)}{b}$, so $a-c=2a=f-d\implies c-d=a-f\implies CD\parallel AF$, and we're done.

 \textbf{Comment.} It appears that taking $(ABC)$ the unit circle is nicer than, say $e=0$ or $(ACE)$ the unit circle (which may not even be reasonably tractable).

 


	\item (\textit{USAMO 2011}) Consider the assertion that for each positive integer $n \ge 2$, the remainder upon dividing $2^{2^n}$ by $2^n - 1$ is a power of 4.  Either prove the assertion or find (with proof) a counterexample.


 



\textbf{Solution}

We will show that $n = 25$ is a counter-example.

 Since $\textstyle 2^n \equiv 1 \pmod{2^n - 1}$, we see that for any integer $k$, $\textstyle 2^{2^n} \equiv 2^{(2^n - kn)} \pmod{2^n-1}$. Let $0 \le m < n$ be the residue of $2^n \pmod n$. Note that since $\textstyle m < n$ and $\textstyle n \ge 2$, necessarily $\textstyle 2^m < 2^n -1$, and thus the remainder in question is $\textstyle 2^m$. We want to show that $\textstyle 2^m \pmod {2^n-1}$ is an odd power of 2 for some $\textstyle n$, and thus not a power of 4.

 Let $\textstyle n=p^2$ for some odd prime $\textstyle p$. Then $\textstyle \varphi(p^2) = p^2 - p$. Since 2 is co-prime to $\textstyle p^2$, we have
 \[{2^{\varphi(p^2)} \equiv 1 \pmod{p^2}}\]and thus
 \[\textstyle 2^{p^2} \equiv 2^{(p^2 - p) + p} \equiv 2^p \pmod{p^2}.\]

 Therefore, for a counter-example, it suffices that $\textstyle 2^p \pmod{p^2}$ be odd. Choosing $\textstyle p=5$, we have $\textstyle 2^5 = 32 \equiv 7 \pmod{25}$. Therefore, $\textstyle 2^{25} \equiv 7 \pmod{25}$ and thus
 \[\textstyle 2^{2^{25}} \equiv 2^7 \pmod {2^{25} - 1}.\]Since $\textstyle 2^7$ is not a power of 4, we are done.

 



\textbf{Solution 2}

Lemma (useful for all situations): If $x$ and $y$ are positive integers such that $2^x - 1$ divides $2^y - 1$, then $x$ divides $y$.Proof: $2^y \equiv 1 \pmod{2^x - 1}$. Replacing the $1$ with a $2^x$ and dividing out the powers of two should create an easy induction proof which will be left to the reader as an Exercise.

 Consider $n = 25$. We will prove that this case is a counterexample via contradiction.

 Because $4 = 2^2$, we will assume there exists a positive integer $k$ such that $2^{2^n} - 2^{2k}$ divides $2^n - 1$ and $2^{2k} < 2^n - 1$. Dividing the powers of $2$ from LHS gives $2^{2^n - 2k} - 1$ divides $2^n - 1$. Hence, $2^n - 2k$ divides $n$. Because $n = 25$ is odd, $2^{24} - k$ divides $25$. Euler's theorem gives $2^{24} \equiv 2^4 \equiv 16 \pmod{25}$ and so $k \ge 16$. However, $2^{2k} \geq 2^{32} > 2^{25} - 1$, a contradiction. Thus, $n = 25$ is a valid counterexample.

 


	\item (\textit{USAMO 2011}) Let $P$ be a given point inside quadrilateral $ABCD$.  Points $Q_1$ and $Q_2$ are located within $ABCD$ such that $\angle Q_1 BC = \angle ABP$, $\angle Q_1 CB = \angle DCP$, $\angle Q_2 AD = \angle BAP$, $\angle Q_2 DA = \angle CDP$.  Prove that $\overline{Q_1 Q_2} \parallel \overline{BA}$ if and only if $\overline{Q_1 Q_2} \parallel \overline{CD}$.


 



\textbf{Solution 1}

First, if $\overline{AB} \parallel \overline{CD}$, then the proof is trivial. Assume $\overleftrightarrow{AB} \cap \overleftrightarrow{CD} = X$. We will show that $\overleftrightarrow{Q_1Q_2}$ passes through $X$, so $\overline{Q_1Q_2} \not \parallel \overline{AB}, \overline{CD}$. Note that $Q_1$ is the isogonal conjugate of $P$ with respect to $\triangle BCX$. Similarly, $Q_2$ is the isogonal conjugate of $P$ with respect to $\triangle ADX$. Thus, $Q_1, Q_2, X$ lie on the reflection of ${XP}$ over the angle bisector of $\angle X$, as desired.

 -sillybone

 -This proof is extremely similar to the one in Evan Chen's \href{https://web.evanchen.cc/exams/USAMO-2011-notes.pdf}{https://web.evanchen.cc/exams/USAMO-2011-notes.pdf}. Please don't sue me.

 



\textbf{Solution 2}

Lemma. If $AB$ and $CD$ are not parallel, then $AB, CD, Q_1 Q_2$ are concurrent.

 Proof. Let $AB$ and $CD$ meet at $R$. Notice that with respect to triangle $ADR$, $P$ and $Q_2$ are isogonal conjugates (this can be proven by dropping altitudes from $Q_2$ to $AB$, $CD$, and $AD$ or $BC$ depending on where $R$ is). With respect to triangle $BCR$, $P$ and $Q_1$ are isogonal conjugates. Therefore, $Q_1$ and $Q_2$ lie on the reflection of $RP$ in the angle bisector of $\angle{DRA}$, so $R, Q_1, Q_2$ are collinear. Hence, $AB, CD, Q_1 Q_2$ are concurrent at $R$.

 Now suppose $Q_1 Q_2 \parallel AB$ but $Q_1 Q_2$ is not parallel to $CD$. Then $AB$ and $CD$ are not parallel and thus intersect at a point $R$. But then $Q_1 Q_2$ also passes through $R$, contradicting $Q_1 Q_2 \parallel AB$. A similar contradiction occurs if $Q_1 Q_2 \parallel CD$ but $Q_1 Q_2$ is not parallel to $AB$, so we can conclude that $Q_1 Q_2 \parallel AB$ if and only if $Q_1 Q_2 \parallel CD$.

 



\textbf{Solution 3}

First note that $\overline{Q_1 Q_2} \parallel \overline{AB}$ if and only if the altitudes from $Q_1$ and $Q_2$ to $\overline{AB}$ are the same, or $|Q_1B|\sin \angle ABQ_1 =|Q_2A|\sin \angle BAQ_2$.  Similarly $\overline{Q_1 Q_2} \parallel \overline{CD}$ iff $|Q_1C|\sin \angle DCQ_1 =|Q_2D|\sin \angle CDQ_2$.

 
If we define $S =\frac{|Q_1B|\sin \angle ABQ_1}{|Q_2A|\sin \angle BAQ_2} \times \frac{|Q_2D|\sin \angle CDQ_2}{|Q_1C|\sin \angle DCQ_1}$, then we are done if we can show that S=1.

 
By the law of sines, $\frac{|Q_1B|}{|Q_1C|}=\frac{\sin\angle Q_1CB}{\sin\angle Q_1BC}$ and $\frac{|Q_2D|}{|Q_2A|}=\frac{\sin\angle Q_2AD}{\sin\angle Q_2DA}$.

 
So, $S=\frac{\sin \angle ABQ_1}{\sin \angle BAQ_2}\cdot\frac{\sin \angle CDQ_2}{\sin \angle DCQ_1}\cdot\frac{\sin \angle BCQ_1}{\sin \angle CBQ_1}\cdot\frac{\sin \angle DAQ_2}{\sin \angle ADQ_2}$

 
By the terms of the problem, $S=\frac{\sin \angle PBC}{\sin \angle PAD}\cdot\frac{\sin \angle PDA}{\sin \angle PCB}\cdot\frac{\sin \angle PCD}{\sin \angle PBA}\cdot\frac{\sin \angle PAB}{\sin \angle PDC}$.  (If two subangles of an angle of the quadrilateral are equal, then their complements at that quadrilateral angle are equal as well.)   

 
Rearranging yields $S= \frac{\sin \angle PBC}{\sin \angle PCB}\cdot\frac{\sin \angle PDA}{\sin \angle PAD}\cdot\frac{\sin \angle PCD}{\sin \angle PDC}\cdot\frac{\sin \angle PAB}{\sin \angle PBA}$.

 
Applying the law of sines to the triangles with vertices at P yields $S=\frac{|PC|}{|PB|}\frac{|PA|}{|PD|}\frac{|PD|}{|PC|}\frac{|PB|}{|PA|}=1$.

 



\textbf{Solution 4}

<i>\textbf{Case 1}</i> The lines $AB$ and $CD$ are not parallel. Denote $E = AB \cap CD.$

 
 \[\angle Q_1 BC = \angle ABP, \angle Q_1 CB = \angle DCP \implies\]Point $Q_1$ isogonal conjugate of a point $P$ with respect to a triangle $\triangle EBC \implies$

 $EP$ and $EQ_1$ are isogonals with respect to $\angle BEC.$

 Similarly point $Q_2$ isogonal conjugate of a point $P$ with respect to a triangle $\triangle EAD \implies$

 $EP$ and $EQ_2$ are isogonals with respect to $\angle BEC.$

 Therefore points $E, Q_1, Q_2$ lies on the isogonal $EP$ with respect to $\angle BEC \implies$

 $Q_1Q_2$ is not parallel to $AB$ or $CD.$

 <i>\textbf{Case 2}</i> $AB||CD.$ We use <i>\textbf{The isogonal theorem in case parallel lines}</i> and get 
 \[AB||Q_1Q_2 || CD. \blacksquare\]

 \textbf{vladimir.shelomovskii@gmail.com, vvsss}

 


	\item (\textit{USAMO 2011}) Let $A$ be a set with $|A| = 225$, meaning that $A$ has 225 elements.  Suppose further that there are eleven subsets $A_1$, $\dots$, $A_{11}$ of $A$ such that $|A_i | = 45$ for $1 \le i \le 11$ and $|A_i \cap A_j| = 9$ for $1 \le i < j \le 11$.  Prove that $|A_1 \cup A_2 \cup \dots \cup A_{11}| \ge 165$, and give an example for which equality holds.


 



\textbf{Solution}

\textbf{Existence} Note that $\textstyle \binom{11}3 = 165,$ and so it is natural to consider placing one element in each intersection of three of the 11 sets. Since each pair of sets is in 9 3-way intersections—one with each of the 9 remaining sets—any two sets will have 9 elements in common. Since $\textstyle \binom{10}2 = 45,$ each set is in 45 triples and thus will have 45 elements. We can now throw in 60 more elements outside the union of the $A_i$ and we are done.

 \textbf{Minimality} As in the proof of PIE, let $I$ be a finite set. Let $\textstyle U = \bigcup_{i\in I} A_i$. For $i\in I$, let $\chi_i$ be the \textit{characteristic function} of $A_i$, that is, for $\alpha \in U,$
 \[\chi_i(\alpha) = \begin{cases} 1 & \alpha \in A_i \\ 0 & \alpha \not\in A_i \end{cases}\]

 For each $\alpha \in U$ let $\textstyle n_\alpha = \sum_{i \in I} \chi_i(\alpha)$, that is the number of subsets $A_i$ of which $\alpha$ is an element.

 If $S \subset I$, let $\textstyle A_S = \bigcap_{i \in S}A_i$. Then the characteristic function of $A_S$ is $\textstyle \prod_{i \in S} \chi_i$.The number of elements of $A_S$ is simply the sum of its characteristic function over all the elements of $U$:
 \[|A_S| = \sum_{\alpha \in U} \prod_{i \in S} \chi_i(\alpha).\]For $0 \le k \le |I|$, consider the sum $N_k$ of $|A_S|$ over all $S\subset I$ with $|S|= k$. This is:
 \[N_k = \sum_{\substack{S\subset I\\|S| = k}} \sum_{\alpha \in U} \prod_{i \in S} \chi_i(\alpha) = \sum_{\substack{S\subset I\\|S| = k}} |A_S|\]reversing the order summation, as an element $\alpha$ that appears in $n_\alpha$ of the $A_i$, will appear in exactly $\textstyle \binom{n_\alpha}{k}$ intersections of $k$ subsets, we get:
 \[\sum_{\substack{S\subset I\\|S| = k}} |A_S| = \sum_{\alpha \in U} \sum_{\substack{S\subset I\\|S| = k}} \prod_{i \in S} \chi_i(\alpha) = \sum_{\alpha \in U} \binom{n_\alpha}k.\]

 Applying the above with $I = \{1, 2, \ldots, 11\},$ and $k=1$, since each of the 11 $A_i$ has 45 elements, we get:
 \[N_1 = \sum_{\alpha \in U} n_\alpha = 11 \cdot 45,\]and for $k = 2$, since each of the 55 pairs $A_i \cap A_j$ has 9 elements, we get:
 \[N_2 = \sum_{\alpha \in U} \binom{n_\alpha}2 = 9 \cdot 55 = 11 \cdot 45.\]

 Therefore
 \begin{align*} s_1 &= \sum_{\alpha \in U} n_\alpha = N_1 = 11\cdot45 \\ s_2 &= \sum_{\alpha \in U} n_\alpha^2 = 2N_2 + N_1 = 3\cdot11\cdot45. \end{align*}
Let $n = |U|$ be the number of elements of $U$.Since for any set of real numbers the mean value of the squares is greater than or equal to the square of the mean value, we have:
 \[\frac{s_2}n \ge \left(\frac{s_1}n\right)^2 \implies n \ge \frac{s_1^2}{s_2} = \frac{11^2\cdot 45^2}{3\cdot11\cdot45} = \frac{11\cdot45}3 = 165.\]

 Thus $|U| \ge 165$ as required.

 



\textbf{Solution 2}

We will count the number of ordered triples, $(A_i,A_j,x)$, where $1\le i,j\le11$ and $x\in A_i \cap A_j$. We know this is equal to $990=11\cdot10\cdot9$­. We can also find that this is $\sum_{k=1}^{225}b_k^2-b_k$, where $b_k$ is the number of the $11$ subsets the $k^{\text{th}}$ element of $A$ is in. Since $\sum_{k=1}^{225}b_k=45\cdot11=495$, we know $\sum_{k=1}^{225}b_k^2=990+495=1485$. Let $n=|A_1 \cup A_2 \cup \dots \cup A_{11}|$. $n$ is equal to the number of $b_k>0$, where $1\le k\le225$. By the QM-AM inequality, we know $\frac{\sum_{k=1}^{225}b_k^2}n\ge\Bigg(\frac{\sum_{k=1}^{225}b_k}n\Bigg)^2\implies\frac{1485}n\ge\Big(\frac{495}n\Big)^2\implies n\ge165$ and that equality occurs when $b_1=b_2=b_3=\cdots=b_{165}=3,b_{166}=b_{167}=\cdots=b_{225}=0$.$\square$

 Solution by randomdude10807

 



\textbf{Solution 3}

Define $x_i$ to be the amount of elements in $A$ such that the element is contained exactly $i$ times among the $11$ subsets $A_j$. We are trying to prove that $\sum_{i=1}^{11}x_i\ge 165$. Now, note that $\sum_{i=1}^{11}ix_i$ is equivalent to the total amount of not necessarily distinct elements among the subsets $A_j$, thus $\sum_{i=1}^{11}ix_i=11|A_j|=495$. Furthermore, note that $\sum_{i=1}^{11}\binom{i}{2}x_i$ is equivalent to the total amount of not necessarily distinct pairs of subsets such that the subsets share an element. In other words, each subset pair should be counted $9$ times in this sum, and there are $\binom{11}{2}$ total distinct pairs of subsets, so $\sum_{i=1}^{11}\binom{i}{2}x_i=9*\binom{11}{2}=495$. Noting that $2\binom{i}{2}+i=i^2$ for all nonnegative integers, we note

 
 \[\sum_{i=1}^{11}i^2x_i=2(\sum_{i=1}^{11}\binom{i}{2}x_i)+(\sum_{i=1}^{11}ix_i)=495*2+495=495*3\]

 Finally, CS inequality gives us

 
 \[(\sum_{i=1}^{11}i^2x_i)(\sum_{i=1}^{11}x_i)\ge(\sum_{i=1}^{11}ix_i)^2\]
 \[(495*3)(\sum_{i=1}^{11}x_i)\ge(495^2)\]
 \[\boxed{(\sum_{i=1}^{11}x_i)\ge 165}\]as desired. 

 One equality case occurs when $x_i=0$ if $i$ is not $3$, and $x_3=165$. Choose any $165$ elements, and represent each with a unique combination of $3$ of the $11$ subsets ($\binom{11}{3}=165$). It is then obvious that each subset has $45$ distinct elements, because after each subset is chosen, there are $\binom{10}{2}=45$ ways to choose the other two subsets to produce a unique element. Similarly, each pair of subsets shares exactly $9$ elements, because after choosing the two subsets, we can choose one of the remaining $9$ subsets to produce the unique element. Hence, proved. $\blacksquare{}$

 


\end{enumerate}